\subsection{The classical balanced DC-SBM oracle and its sharp risk}
\label{dcsbm:classical-oracle}

The fundamental local testing problem used by
Gao, Ma, Zhang and Zhou~\cite{gmzz2018} fixes the target degree parameter
under both candidate labels. Its exact simple-hypothesis likelihood test
therefore corresponds to \(R_{\theta,*}\) in \cref{sec:oracle}.
The next proposition explicitly connects that reference, the profile
reference and true-group neighbor voting in the balanced planted-partition
model. This is an oracle comparison with its own assumptions:
it uses neither the spectral good event \(\mathcal G\) nor a logarithmic
lower bound on \(np\).

\begin{proposition}[Sharp alignment of the classical DC-SBM references]
\label{dcsbm:classical-bridge}
\label{prop:classical-dc-oracle-bridge}
Fix an arbitrary \(K\ge2\), positive bounds
\(\theta_-,\theta_+\), and \(0<r_-<r_+<1\).
Let \(n=Km\), let every group have \(m\) vertices, and suppose
\[
 \theta_-\le\theta_i\le\theta_+,\qquad
 \sum_{i:z_i=a}\theta_i=m\quad\text{for all }a.
\]
Specialize the block matrix \(B=pC\) to
\begin{equation}
 B_{aa}=p,\qquad B_{ab}=q\ (a\ne b),\qquad
 r_-\le q/p\le r_+ .
 \label{dcsbm:classical-model}
\end{equation}
Assume only, for the present oracle comparison,
\begin{equation}
 d=np\longrightarrow\infty,\qquad np^2\le C_1<\infty.
 \label{dcsbm:classical-oracle-range}
\end{equation}
Define \(R_{\mathrm{vote},*}\) as the average error of the local raw-row
rule which knows the neighbor labels and selects the group with the
largest incident block count. Use any fixed tie ordering.
Then, uniformly over this parameter class,
\begin{equation}
 \begin{aligned}
 R_{\theta,*}
 &\asymp R_{\mathrm{prof},*}
 \asymp R_{\mathrm{vote},*}\\
 &\asymp
 \frac1n\sum_{i=1}^n
 \frac1{\sqrt{np}}
 \exp\left\{-\theta_i\frac nK(\sqrt p-\sqrt q)^2\right\}.
 \end{aligned}
 \label{dcsbm:classical-risk}
\end{equation}
The profile reference is exactly the Bernoulli profile in
\eqref{eq:profile-oracle}, with candidate probabilities capped at \(1/2\).
The same comparison holds if the voting counts are normalized by the
self-excluded neighbor-group sizes. Constants depend only on the displayed
fixed bounds, not on \(n\). In particular,
\eqref{dcsbm:classical-oracle-range} permits \(np=o(\log n)\).
\end{proposition}

\begin{proof}
Fix a target \(i\), its true group \(a\), and a competitor \(b\ne a\).
For this proof write
\[
 N_c=\sum_{\substack{j\ne i\\z_j=c}}A_{ij},\qquad
 M_c=\sum_{\substack{j\ne i\\z_j=c}}\theta_j
     =m-\theta_i\ind{c=a},\qquad
 L=\log(p/q).
\]
Thus \(L\) is bounded above and below by positive constants.
Put
\[
 V_{i,ab}=L(N_a-N_b),\qquad
 J_i=\theta_i m(\sqrt p-\sqrt q)^2 .
\]
The assumptions imply \(J_i\asymp d\), uniformly in \(i\).
They also imply \(p\to0\), so every true candidate edge probability
is eventually less than \(1/2\).

\medskip
\noindent\textbf{The exact Bernoulli count transform.}
Under the actual true-row law, independence gives
\begin{align}
 \Psi_{i,ab}(s)
 &:=\log\E_a e^{-sV_{i,ab}}\nonumber\\
 &=\sum_{\substack{j\ne i\\z_j=a}}
   \log\{1+\theta_i\theta_jp(e^{-sL}-1)\}
   +\sum_{j:z_j=b}
   \log\{1+\theta_i\theta_jq(e^{sL}-1)\}.
 \label{dcsbm:classical-cumulant}
\end{align}
Taylor expansion, uniformly for \(s\in[0,1]\), gives
\begin{equation}
 \Psi_{i,ab}(s)
 =m\theta_i\{p(e^{-sL}-1)+q(e^{sL}-1)\}
       +O(np^2+p).
 \label{dcsbm:classical-cumulant-expansion}
\end{equation}
The same error bound holds for the first two derivatives. Bounded theta
and bounded \(L\) make the sum of quadratic Bernoulli remainders
\(O(np^2)\). The additional \(O(p)\) term accounts for the omitted
\(\theta_i\) in \(M_a=m-\theta_i\).

The leading transform in
\eqref{dcsbm:classical-cumulant-expansion} is minimized at \(s=1/2\),
with minimum \(-J_i\). Its endpoint derivatives have opposite signs,
each of magnitude at least \(cd\), and its second derivative is
between positive constant multiples of \(d\).
Since \(np^2+p=o(d)\), the exact minimizing saddle stays in a fixed
compact subinterval of \((0,1)\). Convexity makes it the global
minimum over \(s\ge0\). In particular,
\begin{equation}
 I^{\mathrm{cnt}}_{i,ab}
 :=-\min_{s\ge0}\Psi_{i,ab}(s)
 =J_i+O(np^2+p)=J_i+O(1).
 \label{dcsbm:classical-information}
\end{equation}
One can also obtain this bound directly by taking minima in the
uniform transform approximation.

Under the exact optimizing tilt, the Bernoulli probabilities on the
two relevant groups are comparable to \(p\). Their sum is \(O(d)\),
and the variance of \(V_{i,ab}\) is comparable to \(d\).
Consequently \cref{lem:sharp} applies to
\eqref{dcsbm:classical-cumulant} without any spectral event. For every
fixed \(x_0<\infty\), uniformly over \(|x|\le x_0\),
\begin{equation}
 c_{x_0}d^{-1/2}e^{-J_i}
 \le \PP_a(V_{i,ab}<x)
 \le \PP_a(V_{i,ab}\le x)
 \le C_{x_0}d^{-1/2}e^{-J_i}.
 \label{dcsbm:classical-count-tail}
\end{equation}
The strict lower bound from that lemma uses an open interval of fixed
length under the tilt. Thus lattice atoms and an arbitrary fixed
tie ordering cause no loss of the lower bound.

\medskip
\noindent\textbf{Comparison with the known-target likelihood.}
Let
\[
 \mathcal T_c=\sum_{j\ne i}\theta_j B_{c,z_j}
              =pT_c^{(i)}.
\]
By the exact group normalization,
\begin{equation}
 \begin{aligned}
 \mathcal T_a&=m[p+(K-1)q]-\theta_ip,\\
 \mathcal T_b&=m[p+(K-1)q]-\theta_iq.
 \end{aligned}
 \label{dcsbm:classical-exposure}
\end{equation}
In particular \(\mathcal T_a,\mathcal T_b\asymp d\),
\(\mathcal T_a-\mathcal T_b=O(p)\), and
\(\log(\mathcal T_a/\mathcal T_b)=O(n^{-1})\).

Fix a sufficiently large constant \(L_0\) to be chosen below, and consider
the single-target event
\[
 \mathcal E_i=\{D_i\le L_0d\}.
\]
The exact known-target Bernoulli contrast is
\[
 L^{ab}_{i,\mathrm{known}}
 =V_{i,ab}+
 \sum_{j\ne i}(1-A_{ij})
 \log\frac{1-\theta_i\theta_jB_{a,z_j}}
          {1-\theta_i\theta_jB_{b,z_j}} .
\]
Expansion of these nonedge logarithms on \(\mathcal E_i\) gives
\begin{align}
 L^{ab}_{i,\mathrm{known}}
 &=V_{i,ab}-\theta_i(\mathcal T_a-\mathcal T_b)
       +O(D_ip+np^2)\nonumber\\
 &=V_{i,ab}+O(np^2+p).
 \label{dcsbm:classical-known-contrast}
\end{align}
The \(O(D_ip)\) term restores the incident-edge terms in the linear
exposure; the quadratic remainder is \(O(np^2)\).

\medskip
\noindent\textbf{Comparison with the capped Bernoulli profile.}
For a candidate \(c\), put \(b_{cj}=\theta_jB_{c,z_j}\).
The admissible interval in \eqref{eq:profile-oracle} is
\[
 0\le t\le t_{\max,c}:=\frac1{2\max_{j\ne i}b_{cj}}.
\]
Because \(b_{cj}\asymp p\), its upper endpoint tends to infinity.
For \(t>0\) the derivative of the exact raw-row likelihood is
\[
 \frac{D_i}{t}
 -
 \sum_{j\ne i}(1-A_{ij})
       \frac{b_{cj}}{1-tb_{cj}} .
\]
On \(\mathcal E_i\), there are at least \(n-1-L_0d\ge n/2\)
nonedges for all sufficiently large \(n\). Their total intensity is
at least \(cd\), uniformly over candidates. Hence the derivative is
negative for every admissible \(t\) above a fixed constant depending
only on \(L_0\) and the structural bounds. Every exact maximizing
\(t\) is therefore bounded. The cap at \(1/2\) is inactive on this
bounded interval for all large \(n\).

Uniformly over the resulting bounded \(t\) interval, the likelihood is
\begin{equation}
 D_i\log t+\sum_{j\ne i}A_{ij}\log b_{cj}
      -t\mathcal T_c+O(np^2).
 \label{dcsbm:classical-profile-expansion}
\end{equation}
Indeed, restoring the edge terms in the linear exposure costs
\(O(D_ip)\), and the sum of quadratic nonedge remainders costs
\(O(np^2)\). Both are \(O(np^2)\) on \(\mathcal E_i\).
The leading optimizer \(D_i/\mathcal T_c\) is bounded as well, so
the same error bound holds after maximization.
Terms \(D_i\log D_i-D_i\) and
\(\sum_{j\ne i}A_{ij}\log\theta_j\) cancel in a contrast.
Thus the exact capped profile satisfies
\begin{align}
 L^{ab}_{i,\mathrm{prof}}
 &=V_{i,ab}
   -D_i\log(\mathcal T_a/\mathcal T_b)+O(np^2)\nonumber\\
 &=V_{i,ab}+O(np^2+p)
 \quad\text{on }\mathcal E_i.
 \label{dcsbm:classical-profile-contrast}
\end{align}
For \(D_i=0\), the raw row is identically zero; \(t=0\) maximizes
every candidate likelihood with value zero. The displayed formulas
then hold with the continuous convention \(0\log0=0\).

\medskip
\noindent\textbf{Single-row exceptions, directional errors, and averaging.}
Every mean degree is at most \(\lambda_+d\) for a fixed \(\lambda_+\).
For any prescribed fixed \(H\), a Bernoulli Chernoff bound permits a
fixed choice of \(L_0\) such that
\begin{equation}
 \sup_i\PP(\mathcal E_i^c)\le e^{-Hd}.
 \label{dcsbm:classical-single-exception}
\end{equation}
Also \(J_i\le C_Jd\) uniformly. Choose \(H>C_J+1\).
It follows from \(d\to\infty\) that
\[
 \sup_i
 \frac{\PP(\mathcal E_i^c)}
      {d^{-1/2}e^{-J_i}}
 \le \sqrt d\,e^{-(H-C_J)d}\longrightarrow0 .
\]
This calculation concerns one target at a time; it introduces no
factor \(n\) and requires no logarithmic lower bound on \(d\).

By \eqref{dcsbm:classical-known-contrast} and
\eqref{dcsbm:classical-profile-contrast}, a fixed constant \(B_0\)
bounds the difference between \(V_{i,ab}\) and either exact contrast
on \(\mathcal E_i\). Thus for either
\(L^{ab}_i=L^{ab}_{i,\mathrm{known}}\) or
\(L^{ab}_i=L^{ab}_{i,\mathrm{prof}}\),
\begin{align*}
 \PP_a(V_{i,ab}<-B_0)-\PP_a(\mathcal E_i^c)
 &\le \PP_a(L^{ab}_i<0)\\
 &\le \PP_a(L^{ab}_i\le0)\\
 &\le \PP_a(V_{i,ab}\le B_0)+\PP_a(\mathcal E_i^c).
\end{align*}
Applying \eqref{dcsbm:classical-count-tail} and then
\eqref{dcsbm:classical-single-exception} proves that both the strict
and nonstrict directional errors of either oracle are comparable
to \(d^{-1/2}e^{-J_i}\).
These are unconditional event inclusions; no independence from
\(\mathcal E_i\) is asserted or needed.

For the raw block-count vote, the pairwise contrast is \(V_{i,ab}\)
itself. If counts are normalized by the self-excluded group sizes,
multiplying all candidate scores by \(mL\) changes the pairwise
contrast by at most \(CLD_i/n=O(p)\) on \(\mathcal E_i\), since
those sizes are \(m\) or \(m-1\). The same tail comparison therefore
applies to either voting convention.

For each of these coherent families of candidate scores, a strict
pairwise defeat of the true group implies an error, and an error
implies a nonstrict defeat by some competitor. Taking the maximum
of the strict pair probabilities for a lower bound and their
nonstrict sum for an upper bound proves that each per-target error
is comparable to \(d^{-1/2}e^{-J_i}\), since \(K\) is fixed.
Averaging these bounds over \(i\) proves
\eqref{dcsbm:classical-risk}.
\end{proof}

\begin{remark}[Oracle alignment and the algorithmic range]
\label{dcsbm:classical-scope}
The degree-dependent exponent in \eqref{dcsbm:classical-risk} is the
balanced DC-SBM exponent appearing in the minimax analysis of
Gao, Ma, Zhang and Zhou~\cite{gmzz2018}. Their exponent-level
minimax expression does not, by itself, give the sharp prefactor
or a constant-ratio theorem for every parameter point.
The proposition supplies the exact local-oracle comparison on the
fixed-separation, bounded-theta class stated here.

For \eqref{dcsbm:classical-model}, the finite-\(n\) exposures in
\eqref{eq:equal-exposure} are exactly equal. Hence, on the overlap
with the algorithmic range of \cref{thm:dc},
\cref{cor:known,dcsbm:classical-bridge} give
\[
 \E\ell(\widehat z_D,z)\le C R_{\theta,*},
 \qquad R_{\theta,*}\asymp R_{\mathrm{prof},*}.
\]
The absence of a logarithmic lower bound in the oracle proposition
does not remove that bound from the spectral algorithm theorem.
Extending the algorithm requires its own sparse spectral and
likelihood-transfer estimates.
\end{remark}
